\documentclass[10pt,a4paper]{amsart}
\usepackage[margin=19mm]{geometry}
\usepackage{amsmath,amssymb,mathtools}
\usepackage[scr=rsfs,cal=euler]{mathalfa}
\usepackage{xfrac}
\usepackage[unicode,hidelinks,bookmarks=false]{hyperref}
\newcommand{\ie}{i.e.\ }
\newcommand{\cf}{cf.\ }
\newcommand{\resp}{respectively\ }
\newcommand{\Subsection}{\subsection}
\providecommand{\nobreakdash}{\nobreak\hbox{-}}
\setlength{\emergencystretch}{2em}
\multlinegap0pt
\usepackage{mathtools}
\usepackage{xcolor}
\newtheorem{lemma}{Lemma}
\newcommand{\R}{{\mathbb R}} \newcommand{\Q}{{\mathbb Q}} \newcommand{\Z}{{\mathbb Z}} \newcommand{\N}{{\mathbb N}}
\title{Extreme and exposed points of shift-invariant spaces generated by Gaussian kernel and hyperbolic secant}
\author{Markus Val{\aa}s Hagen, Alexander Ulanovskii, Denis Zelent, Ilya Zlotnikov}
\date{Source: arXiv:2603.04268v1 (CC BY 4.0)}
\begin{document}
\maketitle

\noindent\emph{Source and licence.} Extreme and exposed points of shift-invariant spaces generated by Gaussian kernel and hyperbolic secant. Version arXiv:2603.04268v1 (CC BY 4.0), distributed by arXiv under the Creative Commons Attribution 4.0 International licence, http://creativecommons.org/licenses/by/4.0/, as recorded in the arXiv licence metadata for this exact version and shown on the abstract page https://arxiv.org/abs/2603.04268v1. Full licence notice: \texttt{licenses/arxiv-2603.04268-CC-BY-4.0.txt}.\par\smallskip

\noindent\textit{Editorial scope.} Counted unit: the article's lemma lgf (source0.tex lines 374--376). The article's complete original proof of it is delivered with it (source0.tex lines 378--411, one \texttt{proof} environment of 34 lines). No mathematics is written, repaired or re-derived: the statement and the proof are the article's own lines, in the article's own notation, and no retained line is substituted or re-flowed.  Retained uncounted as the source-local context the proof needs: lines 134--136.  Nothing was omitted from any retained span, so the package carries no omission record.  No retained line contains a citation, so the wrapper carries no bibliography and no bibliography entry.  No retained line contains a figure environment, an image-inclusion command or drawing code, so no image is delivered and the package declares no asset dependency.  Editorial formatting only, in addition: an AMS \texttt{amsart} preamble that restates the article's own declarations and macros used by the retained lines, namely the environments \texttt{lemma} on separate counters, together with the standard packages those lines need.  The retained environments print in this document as Lemma 1 in order of appearance, whereas the article prints the same items with the numbering of its own full document.  The complete native source of the article as distributed in the arXiv e-print for this exact version is bundled unchanged as \texttt{sources/195730/source0.tex}.
\begin{equation}\label{eq:norm_equiv}
    K_1 \|c\|_{l^p} \le \|f\|_{L^p(\R)} \le K_2 \|c\|_{l^p},\quad \text{ for every $f \in V^p_{\Gamma}(g).$}
\end{equation}

\begin{lemma}\label{lgf}
    Assume $f\in V^\infty(G_a)\cap L^1(\R)$. Then $f\in V^1(G_a).$
\end{lemma}

\begin{proof} We set $a=1.$

  Assume $f\in V^\infty(G)\cap L^1(\R)$. Fix any positive $\varepsilon<1$ and consider the function 
 $$f_\epsilon(x):=f\left(\frac{x}{1+\varepsilon}\right)e^{-\varepsilon x^2/(1+\varepsilon)^2}=\sum_{n\in\Z}c_n(f)e^{-(x-n)^2/(1+\varepsilon)-\varepsilon n^2/(1+\varepsilon)}.$$Clearly, $f_\epsilon\in V^1(G_{1/(1+\varepsilon)})$ and $\|f_\varepsilon\|_1\leq (1+\varepsilon)\|f\|_1\leq 2\|f\|_1.$ Note also that the Fourier transform of $f_\varepsilon$ is given by
 $$\hat f_\varepsilon(t)=\int_\R e^{-2\pi i tx}f_\varepsilon(x)\,dx=\sqrt{\pi(1+\varepsilon)}\left(\sum_{n\in\Z}c_n(f)e^{-\varepsilon n^2/(1+\varepsilon)}e^{-2\pi i n t}\right) e^{-\pi^2(1+\varepsilon)t^2}.$$

    Since the trigonometric system $\{e^{2\pi int}\}_{n\in\Z}$ is not complete in $L^2(I),$ for every interval $I$ of length $>1,$ there is a function $\Psi\in C^\infty(\R)$ which vanishes for $|t|\geq 2$ and satisfies 
    $$\int_\R e^{2\pi int}\Psi(t)\,dt=
    \begin{cases}
 \  1,& n=0;\\
 \ 0,& \ n\in\Z\setminus\{0\}.
\end{cases}
   $$

For every $\alpha>0$ we set $$\Psi_\alpha(t):=\sqrt{\frac{\alpha}{\pi}}e^{\pi^2 t^2/\alpha}\Psi(t),$$and denote by $\psi_\alpha$ its inverse Fourier transform. Clearly, $\Psi_\alpha\in C^\infty(\R)$, and   it vanishes for $|t|\geq 2 $, so that $\psi_\alpha$ is a Schwartz function for every $\alpha>0.$ From the Plancherel's theorem, we get
$$\int_\R e^{-\alpha(x-n)^2}\psi_\alpha(x-m)\,dx=\int_\R e^{-\alpha(x-n+m)^2}\psi_\alpha(x)\,dx=\int_\R e^{2\pi i(n-m)t}\Psi(t)\,dt= \begin{cases}
 \  1,& n=m;\\
 \ 0,&  n \neq m.
\end{cases}$$

We now assume that $\alpha=1/(1+\varepsilon)\in[1/2,1]$. We have
\begin{equation}\label{eq:ortck}
    \int_\R f_\varepsilon(x) \psi_{1/(1+\varepsilon)}(x-k) \, dx=c_k(f)e^{-\varepsilon n^2/(1+\varepsilon)}.
\end{equation}

Clearly, 
$$\|\psi_\alpha\|_{L^\infty(\R)}+ \| x^2\psi_\alpha(x)\|_{L^\infty(\R)}\leq \|\Psi_\alpha\|_{L^1(\R)}+ \|(\Psi_\alpha)''\|_{L^1(\R)}<C,$$where $C$ does not depend on $\varepsilon.$
Therefore, 
 \begin{equation}\label{epsi}\sup_{x\in\R}\sum_{n\in\Z}|\psi_{\alpha}(x-n)| <C,\quad \text{for every } \alpha=1/(1+\varepsilon)\in [1/2,1].\end{equation}


Consider the integral $$I_{\varepsilon}:=\int_\R f_\varepsilon(x)\sum_{k\in\Z}\frac{\overline{c_k(f)}}{|c_k(f)|}\psi_{1/(1+\varepsilon)}(x-k)\,dx.$$By \eqref{epsi}, 
\begin{equation}\label{e0e}|I_\varepsilon|<C\|f_\varepsilon\|_1\leq 2C\|f\|_1,\quad 0<\varepsilon\leq 1.\end{equation} On the other hand, since  the sequence $\{c_n(f)\exp(-\varepsilon n^2/(1+\varepsilon))\}\in l^1(\Z)$, by the right-hand-side inequality in \eqref{eq:norm_equiv}, the function  $$\sum_{n\in\Z}|c_n(f)|e^{-\varepsilon n^2/(1+\varepsilon)}e^{-(x-n)^2/(1+\varepsilon)}$$belongs to $L^1(\R).$ By the Dominated convergence theorem, integrating term for term  and using \eqref{eq:ortck}, we obtain $$I_\varepsilon =\sum_{n\in\Z}|c_n(f)|e^{-\varepsilon n^2/(1+\varepsilon)}.$$ Using \eqref{e0e} and  letting $\varepsilon \to 0$, we get $\{c_n(f)\}\in l^1(\Z),$ so that $f\in V^1(G).$
\end{proof}

\end{document}
